\chapter{Motion and Its Records}
\label{chap:motion-and-its-records}

To a finite instrument, motion is not a curve drawn in a pre-existing space; it is a record --- a ledger of
events, each one a count laid down in order. Everything this book says about physics is read off such records
by one apparatus, and this chapter is where that apparatus first goes to work. Before the five effects of
motion, then, it is worth meeting the instrument itself: not as a metaphor, but as a real piece of equipment
with a definite way of reading the world, a definite resolution, and a definite, honest account of what it
can and cannot claim. Every experiment in this book --- through fields, heat, the quantum, relativity, and
the foundations --- is this same instrument at work, and the reader who holds it in mind from the first count
will recognize it behind the last.

The first question any record forces is which of its features are facts about the world and which are facts
about the bookkeeping we lay over it. The answer, stated once here and used everywhere after, is that the
physical content of a record is exactly what survives relabeling: rename the events, shift the origin, boost
to a moving frame, and whatever is left unchanged is the physics. This chapter establishes that discipline through five
effects of motion --- but it begins with the apparatus that will enforce it.

\section{The Instrument}
\label{se:the-instrument}

Picture an instrument on a bench. It does only one primitive thing, and everything else is built from that
single act. The primitive
act is to record a \emph{difference} --- to mark that two configurations are not the same. The first mark it
ever makes is the first distinction it can draw, and from that one recorded difference its entire repertoire
follows. Six features describe how it works, and the whole book is their elaboration.

\emph{It is finite.} It works with a finite number of parts and a finite memory, and every step it takes
touches only what it can actually hold --- never a completed infinity, never an unbounded resource. The
features that follow all operate within this one limit, and so does every wall the book later meets.

\emph{It counts.} Every distinction the instrument draws is added to a growing ledger of events --- a count.
To measure, for this instrument, is to increment that count: an outcome exists for it exactly when it adds a
unit to the record. The count is the instrument's only currency, and the book's recurring claim is that an
entire physics can be read off it.

\emph{It has a floor.} The instrument's resolution is finite: there is a smallest difference it can record, a
grid below which two states are, to it, the same state. Nothing it reads is finer than that floor. The
continuous quantities of textbook physics --- a smooth position, an exact real number --- are never what the
instrument holds; they are the limit its counts approach as the grid is imagined refined, a limit it
describes but never reaches.

\emph{It reconciles at boundaries.} When two parts of the record meet and disagree --- when two sub-ledgers
carry incompatible entries at a shared edge --- the instrument does not scan their interiors to adjudicate.
It performs a reconciliation at the boundary, a pivot that fixes a consistent order from the edge alone. Much
of what looks mysterious in the later parts --- a measurement that ``collapses,'' two distant records that
agree without a signal passing between them --- is this boundary reconciliation, seen from outside.

\emph{It carries holonomy.} Send a signal around a closed loop in the record and back to its start, and it
need not return unchanged: it can come back with a sign, a $\pm 1$ that measures how the loop is wound. Which
of the two signs that $\pm 1$ shows is fixed only once an orientation is given to the loop --- traverse it the
other way and the sign flips --- but \emph{that} it returns signed at all, rather than unchanged, is the fact
and not the convention. This residue is not an approximation the instrument makes; in the right experiments it
is a fact the instrument genuinely \emph{proves} --- a discrete, checkable property of the loop. The echo in a maze, the bright peak
of a diffracted electron, the handedness of a fermion, and the detection of a positron are all this same loop
residue, proved.

\emph{It is honest.} For every reading, the instrument states exactly what kind of claim it is making, and it
recognizes only four kinds --- four \emph{ceilings}, four honest floors beneath which it will not pretend to
reach. A reading may be a \emph{finite count obligation}: a balance, a conservation, a closure the record
genuinely owes. It may be a \emph{smooth-shadow analogy}: a continuum law named as the smooth shadow the
discrete count casts as the grid refines, the bridge to the equation called an analogy and never a
derivation. It may be a \emph{no-go}: a wall, a refinement the law forbids. Or it may be a \emph{name}: a
phenomenon written out in full as physics, with the instrument claiming nothing but the label. The instrument
never claims more than the floor it stands on, and naming the floor is half of every reading in this book.

That is the whole apparatus. The five effects of this chapter are its first motions --- the instrument
learning to tell a label from a fact, what it has imposed from what it has found --- and every later chapter is
the same instrument reading a new corner of physics off the same finite count.

\section{The Instrument at Work}
\label{se:instrument-at-work}

Before the named effects, watch the instrument take a single reading, so the abstractions above become a
thing on a bench. Set a ball rolling along a grooved track, and place beside the track a ruler marked in
millimetres and a clock ticking at a fixed rate. The instrument does not see a ball gliding through space; it
sees a sequence of coincidences --- the ball's edge level with \emph{this} mark as the clock shows
\emph{that} tick --- and it records each coincidence as one event, a pair of counts: which ruler-cell, which
clock-tick. A motion, to the instrument, is the resulting ledger,
\begin{equation}
  R = \bigl( (n_1, k_1),\ (n_2, k_2),\ \dots,\ (n_N, k_N) \bigr),
\end{equation}
of $N$ events, each a cell against the ruler paired with a tick against the clock. Nothing continuous has
been recorded; what the instrument holds is a finite table of coincidences, and the notation is a description
of that table, not a function to be evaluated.

Now exercise the instrument's discipline on this one reading. Slide the ruler's zero to a new mark --- the
zero was ours to place, never the world's: every $n_i$ changes, the ledger looks entirely different, and yet
the gaps $n_{i+1} - n_i$ are untouched --- the instrument has relabeled, and the part of the reading that
survives the relabeling is the motion. Coarsen the
ruler, keeping only every other cell: events that fell in the same coarse cell collapse to one, and below the
new spacing the instrument can no longer tell two positions apart --- its floor has risen, and what it reads
is honestly blunter. Set a second clock drifting uniformly against the first and re-file the events by it:
every tick-count shifts, but the differences of differences --- the accelerations --- do not move, because a
uniform drift is precisely the boost the instrument's invariance class is built to ignore. And bring a second
ball's ledger alongside the first, each with its own ruler and clock, and ask for their relative motion: the
instrument reconciles the two tables at their shared clock and keeps only the difference of their
position-counts.

Everything the rest of the chapter names --- the Aristotle invariance, the Galilean class, the Cartesian cell
and pair, the relational velocity --- has just happened, once, on this bench. The five effects that follow
are not five new machines; they are five things to notice about the one reading the instrument has already
taken.

\section{The Aristotle Effect: relabeling is not a fact}
\label{se:aristotle}

The oldest confusion about motion is to mistake the description for the thing described, and it is the first
confusion the instrument is built to resist. Aristotle's effect, in the form a record forces, is the
recognition that re-enumerating the events of a motion --- an act of ours, not of the motion --- changes the
labels and nothing else. Let a record be
an ordered list of events $R = (e_1, e_2, \dots, e_N)$, and let a relabeling be a permutation $\sigma$ acting
on the indices, $\sigma \cdot R = (e_{\sigma(1)}, \dots, e_{\sigma(N)})$. The notation is descriptive:
$\sigma$ is the \emph{name} of an act of re-enumeration, and the question it poses is which of the record's
features are blind to it. A quantity is a \emph{fact} about the motion only if it is unchanged by every
admissible $\sigma$,
\begin{equation}
  Q(\sigma \cdot R) = Q(R) \qquad \text{for all admissible } \sigma .
\end{equation}
The instrument's own count is the first such invariant. Writing $N(R)$ for the number of recorded events,
\begin{equation}
  N(\sigma \cdot R) = N(R),
\end{equation}
because a permutation moves entries without creating or destroying them --- it shuffles the ledger, and the
height of the ledger does not care how it is shuffled. More generally, any functional that depends only on
the multiset of events --- on \emph{which} events occurred, not on the order in which they were filed --- is
invariant; any functional that depends on the filing order is not. The equation computes nothing; it draws a
line, between the part of the record that is the instrument's bookkeeping and the part that is the world.

What the finite experiment actually secures here is modest and exact, and naming its honest floor is the
point: it is a finite count obligation. The instrument does not prove a theorem about a continuum of
relabelings; it exhibits a count that is preserved when the record is re-enumerated, and a count is a finite,
checkable thing --- the instrument can simply re-shuffle its ledger and confirm the height is unchanged. The
reading is correspondingly sharp. What is invariant under coordinated re-enumeration constitutes empirical
fact; what an enumeration introduces, an enumeration may remove, and so cannot be a fact about the motion.
The labels are ours; the count is the world's. This is the instrument's founding discipline --- and the
book's final wall, that no renaming of the record can raise its entropy, is this same sentence met again at
the very end, this time as a thing that cannot be done.

\section{The Galileo Effect: objectivity is an invariance class}
\label{se:galileo}

Galileo sharpened the same recognition into a principle of relativity, and in doing so handed the instrument
its first invariance class. The admissible relabelings of a record include not only re-enumerations but
changes of the frame from which the record is taken --- a shift of origin, a rotation of axes, and,
decisively, a boost to a uniformly moving frame. The Galilean transformation between two frames in relative
motion at velocity $v$ is written
\begin{equation}
  x' = x - v t, \qquad t' = t .
\end{equation}
It is essential to read this equation the way the instrument does. It is not a recipe to be differentiated
for a velocity rule; it is the \emph{name} of a relabeling --- a description of how one frame's ledger is
rewritten as another's --- and what matters is not any number ground out of it but what it leaves
\emph{fixed}. A uniform boost slides every recorded position by the same drifting amount, so it cannot touch
any feature of the record built from differences of differences: the spacing between successive positions
shifts by a constant, and the change in that spacing does not change at all. The acceleration is therefore
blind to the boost, and Newton's law, written in the acceleration,
\begin{equation}
  F = m\,\ddot{x},
\end{equation}
reads the same in the boosted frame as in the original --- not because we solved for $F'$ and found it equal,
but because the law is assembled from exactly the part of the record the boost leaves alone. Gather the full
set of admissible changes into a group $\mathcal{G}$ acting on records; then a statement $S$ is
\emph{physical} precisely when it is a fixed point of the action,
\begin{equation}
  S(g \cdot R) = S(R) \qquad \text{for all } g \in \mathcal{G}.
\end{equation}
This fixed-point equation is the whole content of Galilean objectivity, stated as a picture rather than a
calculation: the physics is the part of the record that holds still while the group moves it.

Here too the honest floor is a finite count obligation: the instrument certifies that a recorded statement
survives a relabeling by checking that the counts it depends on are unchanged, not by quantifying over a
continuum of frames. The reading is Galileo's, made operational. Objectivity is not a property an isolated
statement possesses; it is membership in the invariance class of statements that agree under every admissible
relabeling. What cannot survive relabeling is not thereby false --- it is simply not a fact about the motion,
having been a fact about the frame we read it from. The instrument, which carries the group action in its very
construction,
reads objectivity directly: a statement is objective when the apparatus can relabel its frame and find the
statement unmoved.

\section{The Descartes Effect, I: position quantized against axes}
\label{se:descartes-i}

Descartes gave motion coordinates, and a finite instrument gives those coordinates a smallest division ---
its resolution floor, the first of the instrument's limits to appear. Position is recorded not as a real
number but as a cell index against a fixed axis of spacing $\Delta$,
\begin{equation}
  x \;\longmapsto\; n = \left\lfloor \frac{x}{\Delta} \right\rfloor, \qquad x \in [\,n\Delta,\,(n+1)\Delta\,),
\end{equation}
so that each recorded event is a quantized instance of position, accurate to within a cell. The floor
brackets are descriptive, not computational: they say \emph{which cell} the position falls in, and they
deliberately forget where in the cell. Below $\Delta$ the instrument records nothing --- two positions in the
same cell are, to it, one position --- and this is not a defect to be corrected but the instrument's honest
resolution, the very floor that will later forbid an arbitrarily sharp simultaneous reading and set the
zero-point of an oscillator. What is invariant across the quantization is not the cell indices, which depend
on the origin and spacing we lay down, but the geometric relation among the recorded positions --- the trajectory
itself --- which is unchanged under a shift of origin, a rotation of axes, and a uniform rescaling,
\begin{equation}
  x \mapsto R\,x + b, \qquad x \mapsto \lambda\, x .
\end{equation}
The quantized ledger entries vary with the chart; the structural curve they sample remains recoverable from
any of them, exactly as the Galilean section requires.

The finite floor is, again, a count obligation: the instrument secures a quantized ledger of positions whose
geometric relations are preserved, not a continuum of trajectories. And the reading is the bridge the rest of
mechanics will rest on without further comment: a continuous motion is the smooth shadow the quantized ledger
casts as the division is imagined refined,
\begin{equation}
  \gamma(t) \;=\; \lim_{\Delta \to 0}\; n(t)\,\Delta .
\end{equation}
This limit is the instrument's second kind of honest floor making its first appearance --- a smooth-shadow
analogy, named here so it can be relied upon later. The curve is what the cells approach; it is never a cell,
and the instrument that reads the cells describes the curve without ever holding it.

\section{The Descartes Effect, II: independent rulers and the ordered pair}
\label{se:descartes-ii}

A single axis counts along one direction; physical space needs more than one, and the way the counts combine
is itself a physical fact the instrument must respect. Let two independent rulers generate their own stable
counts $N_a$ and $N_b$. The naive move is to add them into a single total $N_a + N_b$, but that flattens the
record: a total cannot say how much came from each ruler, and an instrument that kept only the sum would have
thrown away a distinction it had paid to record. The lawful move is to keep them as an ordered pair,
\begin{equation}
  (N_a, N_b) \in \mathbb{Z}^2,
\end{equation}
which preserves exactly what the flat sum conflates. The order matters, and it matters for a reason that will
grow into one of the book's central themes: advancing along $a$ and then along $b$ is not, in general, the
same act as advancing along $b$ and then along $a$. The two compositions trace different paths, and a record
that keeps the pair remembers the difference while a record that keeps the sum forgets it. That failure of
two operations to commute --- here a harmless bookkeeping fact about which ruler moved first --- is the seed
of holonomy: when the two operations are transports around a loop rather than steps along an axis, their
failure to commute is exactly the $\pm 1$ the instrument will later prove a charged loop to carry. Two
dimensions of geometry are the price of refusing to collapse two independent counts into one.

The honest floor is a finite count obligation --- a pair of counts that must each be recoverable from the
record. The reading is the one Descartes did not state but the record forces: geometry is not assumed by the
theory and then confirmed; it is \emph{forced} by the demand that two independent counts reconcile without
losing which is which. The plane is what you get when you insist on keeping both rulers, and the order in
which the rulers are read is the first faint sign of the loops to come.

\section{The Velocity Effect: velocity as a relational count}
\label{se:velocity}

With positions counted and frames understood, velocity follows --- but not as a substance a body carries.
Velocity is a relational count: the rate at which one record's events accumulate against another's. For a
single body against a clock,
\begin{equation}
  v \;=\; \frac{\Delta x}{\Delta t} \;=\; \frac{x(t_2) - x(t_1)}{t_2 - t_1},
\end{equation}
a difference of position counts over a difference of time counts. The ratio is descriptive: it names how two
of the instrument's tallies --- positions and clock-ticks --- accumulate against each other, and it is
meaningful only once both tallies are entries in a common ledger. Relative velocity is the same construction
applied to two bodies once their records have been merged, and the merging is the instrument's boundary
reconciliation in its simplest form: two separate ledgers are brought to a shared edge and made consistent,
and of the reconciled counts only the \emph{differences} survive,
\begin{equation}
  v_{\mathrm{rel}} \;=\; v_1 - v_2 \;=\; \frac{\Delta(N_1 - N_2)}{\Delta t},
\end{equation}
so that a common drift shared by both bodies cancels and contributes nothing. This is the Galilean invariance
of the earlier section seen from the side of the bodies rather than the frame: the boost that left the law
fixed is here the common drift the difference discards.

The finite floor is a count obligation: the instrument secures a difference of event counts, not a velocity
field on a continuum. And the reading closes the chapter where it opened. A velocity is not a state the world
stores and the instrument reads off; it is a difference of counts the instrument forms after two records are
reconciled at their shared edge, and only differences survive that reconciliation. Motion, recorded honestly,
is relational from the first count to the last.

\bigskip

The instrument is now in view, and so is its discipline. It begins from a single recorded difference; it lays
down a count; it reads motion off that count by keeping what survives relabeling and discarding what does
not; and already, in these five effects of motion, it has shown the first signs of everything to come --- a
resolution floor, a continuum that is only a shadow, a pair of counts whose order will grow into a loop, a
reconciliation at a boundary, and four honest ceilings it will never overclaim. The chapters that follow are
this same apparatus carried outward: to the fields, where the loop residue becomes a holonomy it proves; to
heat, where the count becomes a temperature and the second law a price; to the quantum, where the count
becomes an amplitude and a single particle is proved at a detector; to relativity, where the geometry of the
count bends; and to the foundations, where the instrument is turned on itself and asked what it cannot
escape. Hold the instrument in mind. Everything after this is watching it work.
